\documentclass[conference]{IEEEtran}
\IEEEoverridecommandlockouts
\usepackage[utf8]{inputenc}
\usepackage{cite}
\usepackage{amsmath,amssymb,amsfonts}
\usepackage{graphicx}
\usepackage{booktabs}
\usepackage{url}
\usepackage{balance}

\begin{document}

\title{Why Observational AI-Agent Pull Request Datasets Cannot Support Behavioral Inference: A Critical Analysis of Methodology Failures in Software Engineering Research}

\author{\IEEEauthorblockN{Ahmed Mursal}
\IEEEauthorblockA{\textit{School of Computing} \\
\textit{Edinburgh Napier University}\\
Edinburgh, Scotland \\
40646515@live.napier.ac.uk}
}

\maketitle

\begin{abstract}
We examine the suitability of observational pull request datasets for AI agent behavioral research using the AIDev dataset as a case study. Our analysis reveals fundamental validity failures that preclude meaningful inference: agent attribution exhibits obvious errors (GitHub Copilot: 379 users vs. millions in reality), statistical tests violate independence assumptions due to clustering, and measurement approaches show systematic bias. We demonstrate why large-scale observational studies of AI tools cannot support behavioral claims without controlled experimental designs. This methodology critique provides guidelines for valid AI tool evaluation and warns against common pitfalls in software engineering research that mistake dataset size for validity.
\end{abstract}

\begin{IEEEkeywords}
methodology, software engineering research, AI tools, empirical evaluation, validity threats
\end{IEEEkeywords}

\section{Introduction}

The software engineering research community increasingly relies on large observational datasets to study AI coding tools. The AIDev dataset exemplifies this trend, containing 932,791 pull requests with agent attribution labels that researchers have used to draw behavioral conclusions about different AI tools. However, such observational approaches face fundamental methodological challenges that compromise validity.

This paper examines whether pull request datasets can support reliable behavioral inferences about AI tools. Through systematic analysis of the AIDev dataset, we identify critical failures in: (1) agent attribution reliability, (2) statistical assumption compliance, and (3) measurement validity. These issues collectively demonstrate why observational PR data cannot support behavioral claims about AI tools.

\subsection{The Appeal of Large Datasets}

Large observational datasets offer compelling advantages: massive sample sizes, real-world authenticity, and comprehensive coverage. The AIDev dataset's 932,791 observations across 72,189 developers appears to provide unprecedented statistical power for understanding AI tool behavior. This scale creates an illusion of methodological rigor that can mask fundamental validity threats.

\subsection{Research Questions}

We address three methodological research questions:

\textbf{RQ1: What validity threats prevent reliable inference about AI agent behavior from observational PR datasets?} We systematically identify attribution failures, statistical violations, and measurement biases.

\textbf{RQ2: How do these validity threats manifest in practice using the AIDev dataset?} We provide concrete examples of each threat category with specific evidence.

\textbf{RQ3: What methodological requirements are necessary for valid AI tool behavioral research?} We specify design principles for reliable causal inference.

\section{Methodology Failures in Observational AI Research}

\subsection{RQ1: Systematic Validity Threats}

\subsubsection{Agent Attribution Unreliability}

Observational datasets typically derive agent labels through metadata extraction, commit message parsing, and self-reports. This approach suffers from multiple failure modes:

\textbf{Detection Incompleteness.} Automated attribution can only identify explicitly declared tool usage, missing silent assistance or tools used without obvious signatures.

\textbf{Temporal Inconsistency.} Tool adoption changes rapidly, but retrospective labeling may use contemporary usage patterns to infer historical tool use.

\textbf{Mixed Authorship.} Real development involves human editing of AI-generated code, making pure agent attribution impossible in observational settings.

\textbf{Self-Report Bias.} Developers may over-report or under-report tool usage based on perceived social desirability or organizational policies.

\subsubsection{Statistical Independence Violations}

Standard statistical tests assume independent observations, but PR datasets exhibit multiple clustering structures:

\textbf{Repository Clustering.} PRs within repositories share coding standards, review processes, and maintainer preferences. These systematic similarities violate independence.

\textbf{Developer Clustering.} Multiple PRs from the same developer reflect individual coding style, experience level, and tool preferences rather than independent observations.

\textbf{Temporal Clustering.} PRs within development cycles often follow sequential phases (feature development, testing, documentation) that create dependencies.

\textbf{Project Clustering.} PRs within the same project share requirements, constraints, and quality standards that influence testing practices.

These clustering violations invalidate chi-square tests, t-tests, and other standard procedures that assume independence, leading to inflated significance and misleading confidence intervals.

\subsubsection{Measurement Validity Problems}

Behavioral inference requires valid measurement instruments, but observational approaches suffer systematic bias:

\textbf{Language Ecosystem Bias.} Test detection methods favor popular languages and frameworks while systematically missing others, creating artificial behavioral differences.

\textbf{Context Confusion.} Mixing file paths, commit messages, PR descriptions, and code content creates inconsistent measurement criteria.

\textbf{Quality vs. Quantity.} Observational methods can only measure test presence, not test quality, coverage, or effectiveness.

\textbf{Validation Inadequacy.} Manual validation samples (typically 200-500 observations) cannot detect systematic biases across hundreds of thousands of observations spanning multiple domains.

\subsection{RQ2: Evidence from the AIDev Dataset}

\subsubsection{Attribution Failure Examples}

Examination of the AIDev dataset reveals obvious attribution errors that demonstrate label unreliability:

\textbf{Impossible User Counts.} GitHub Copilot appears with only 379 users despite having millions of actual subscribers. This 1000:1 under-detection indicates severe attribution failure.

\textbf{Implausible Market Share.} Lesser-known tools show more users than established market leaders, contradicting all external adoption data.

\textbf{Extreme Single-User Cases.} Devin shows 29,744 pull requests from a single user, suggesting bulk operations misattributed as agent behavior.

\textbf{Temporal Inconsistencies.} Tools show activity patterns that contradict known release timelines and API availability.

These examples demonstrate that agent labels cannot support group comparisons or behavioral inference.

\subsubsection{Statistical Assumption Violations}

The AIDev dataset exhibits severe clustering that invalidates standard statistical analysis:

\textbf{Repository Concentration.} Major agents are concentrated in small numbers of repositories, violating distributional assumptions.

\textbf{Developer Concentration.} Heavy users contribute hundreds of PRs each, creating extreme non-independence.

\textbf{Project Dependencies.} Sequential PRs within projects follow development workflows that create systematic patterns.

Standard chi-square analysis of such data produces meaningless p-values and inflated effect sizes.

\subsubsection{Measurement Bias Evidence}

Test detection approaches show systematic bias that undermines behavioral claims:

\textbf{Keyword Bias.} Detection favors Python/JavaScript terminology while missing Go, Rust, C++, and domain-specific testing approaches.

\textbf{False Positive Patterns.} Phrases like "tested in production," "integration challenges," and "test data" trigger detection without indicating actual tests.

\textbf{Ecosystem Blindness.} Languages with different testing conventions (property-based testing, contract testing, specification-driven development) are systematically under-detected.

\subsection{RQ3: Requirements for Valid AI Tool Research}

\subsubsection{Experimental Design Requirements}

Valid behavioral research about AI tools requires controlled experimental conditions:

\textbf{Randomized Assignment.} Participants must be randomly assigned to tools to eliminate selection bias and confounding.

\textbf{Controlled Tasks.} Standardized programming tasks ensure fair comparison across conditions.

\textbf{Verified Attribution.} Direct tool instrumentation provides accurate usage measurement without inference.

\textbf{Appropriate Controls.} Baseline conditions without AI assistance enable causal attribution.

\subsubsection{Measurement Best Practices}

\textbf{Language-Stratified Validation.} Test detection must be validated within language ecosystems to avoid systematic bias.

\textbf{Expert Review.} Stratified manual validation across tools, languages, and project types with domain expertise.

\textbf{Process Measurement.} Focus on development processes and intermediate artifacts, not just final outcomes.

\textbf{Quality Assessment.} Evaluate test effectiveness through coverage analysis, mutation testing, or expert review.

\subsubsection{Statistical Approaches}

\textbf{Cluster-Robust Methods.} Use mixed-effects models, cluster-robust standard errors, or hierarchical approaches that account for non-independence.

\textbf{Matched Sampling.} Control confounders through propensity score matching, stratification, or covariate adjustment.

\textbf{Effect Size Focus.} Emphasize practical significance and confidence intervals over p-values.

\textbf{Sensitivity Analysis.} Test robustness to measurement assumptions and modeling choices.

\section{Implications for Software Engineering Research}

\subsection{The Illusion of Big Data Validity}

Large sample sizes create false confidence in observational studies. The AIDev dataset's 932,791 observations appear to provide unprecedented statistical power, but scale cannot overcome fundamental validity threats. In fact, larger samples can amplify systematic biases and create stronger false signals.

\textbf{Sample Size ≠ Validity.} Measurement error, attribution failures, and assumption violations affect large datasets just as severely as small ones. Scale may actually worsen these problems by making manual validation infeasible.

\textbf{Significance Inflation.} Massive samples make even tiny effects statistically significant, misleading researchers about practical importance.

\textbf{Bias Amplification.} Systematic errors compound with sample size, creating stronger false signals that appear more convincing.

\subsection{Common Methodological Pitfalls}

Software engineering researchers frequently fall into predictable traps when analyzing observational AI tool data:

\textbf{Attribution Confidence.} Trusting metadata-derived labels without validation leads to meaningless group comparisons.

\textbf{Statistical Misapplication.} Applying independence-assuming tests to clustered data produces invalid inference.

\textbf{Causal Overreach.} Interpreting correlational patterns as behavioral evidence without controlling for confounders.

\textbf{Validation Theater.} Using small manual validation samples to justify conclusions about massive datasets.

\subsection{Guidelines for Valid Research}

Before conducting observational AI tool research, researchers should:

\subsubsection{Attribution Validation}
\begin{enumerate}
\item Verify tool labels match known adoption patterns
\item Investigate extreme distributions or single-user anomalies  
\item Cross-validate against independent data sources
\item Report attribution uncertainty and sensitivity analyses
\end{enumerate}

\subsubsection{Statistical Diagnosis}
\begin{enumerate}
\item Quantify clustering through intracluster correlation
\item Test independence assumptions using diagnostic statistics
\item Use appropriate methods for clustered data
\item Report effective sample sizes, not just raw counts
\end{enumerate}

\subsubsection{Measurement Validation}
\begin{enumerate}
\item Stratify validation across relevant dimensions
\item Test for systematic bias in detection methods
\item Use domain experts for quality assessment
\item Report measurement uncertainty and limitations
\end{enumerate}

\section{Recommendations for Future Research}

\subsection{Methodological Priorities}

\textbf{Controlled Experiments.} The research community should prioritize randomized controlled trials over observational studies for behavioral claims about AI tools.

\textbf{Within-Subject Designs.} Longitudinal studies tracking individual developers across different tools can control for person-level confounders.

\textbf{Mixed Methods.} Combine quantitative measurement with qualitative investigation of decision-making processes.

\textbf{Replication.} Independent replication of tool evaluation findings across different contexts and populations.

\subsection{Dataset Development}

\textbf{Verified Attribution.} Future datasets should include direct tool instrumentation rather than inferential labeling.

\textbf{Rich Context.} Include repository characteristics, developer experience, project requirements, and organizational factors.

\textbf{Quality Metrics.} Extend beyond presence/absence to measure code quality, test effectiveness, and maintenance outcomes.

\textbf{Longitudinal Structure.} Track changes over time rather than cross-sectional snapshots.

\subsection{Community Standards}

The software engineering research community should establish methodological standards for AI tool evaluation:

\textbf{Validity Checklists.} Standardized criteria for evaluating observational AI tool studies.

\textbf{Reporting Guidelines.} Required disclosure of attribution methods, validation procedures, and assumption violations.

\textbf{Peer Review Training.} Education for reviewers on validity threats specific to AI tool research.

\textbf{Negative Result Acceptance.} Venues should welcome methodology critiques and null findings that advance research quality.

\section{Limitations of This Analysis}

This critique itself has limitations that readers should consider:

\textbf{Single Dataset Focus.} We examine only the AIDev dataset; other datasets may exhibit different validity patterns.

\textbf{Methodology Emphasis.} We focus on internal validity rather than practical utility of observational approaches.

\textbf{Alternative Approaches.} We do not explore all possible methods for addressing validity threats in observational settings.

\textbf{Evolving Landscape.} Rapid changes in AI tools may create new challenges not captured in our analysis.

\section{Conclusion}

This analysis demonstrates that observational pull request datasets cannot support reliable behavioral inferences about AI coding tools. The AIDev dataset, despite its impressive scale, suffers from attribution failures, independence violations, and measurement bias that collectively preclude valid behavioral claims.

\textbf{Core Message.} Dataset size cannot substitute for methodological rigor. Valid behavioral research about AI tools requires controlled experimental designs with verified attribution, appropriate statistical methods, and validated measurement instruments.

\textbf{Call to Action.} The software engineering research community should resist the appeal of large convenience datasets and prioritize methodological validity over sample size. Future research should focus on controlled experiments, within-subject comparisons, and mixed-method approaches that can support causal inference.

\textbf{Positive Outlook.} This critique aims to improve research quality, not discourage investigation. By identifying common pitfalls and establishing validity requirements, we can build a more reliable evidence base for AI tool evaluation and software engineering practice.

\section*{Acknowledgments}

We thank H. Li et al. for making the AIDev dataset publicly available, enabling this methodological analysis. We hope this critique contributes to improved research practices in software engineering tool evaluation.

\balance
\begin{thebibliography}{00}
\bibitem{li2024rise} H. Li et al., "The Rise of AI Teammates in Software Engineering (SE 3.0)," arXiv:2409.18925, 2024.
\bibitem{chen2021evaluating} M. Chen et al., "Evaluating Large Language Models Trained on Code," arXiv:2107.03374, 2021.
\bibitem{bareiss2023code} A. Bareiß et al., "Code generation tools in practice: A user study with GitHub Copilot," in Proc. MSR, pp. 285-297, 2023.
\bibitem{bird2011dont} C. Bird et al., "Don't touch my code! Examining the effects of ownership on software quality," FSE 2011.
\bibitem{bacchelli2013expectations} A. Bacchelli and C. Bird, "Expectations, outcomes, and challenges of modern code review," ICSE 2013.
\end{thebibliography}

\end{document}